\section{Mass coordinates and spherical electrovac}
\label{c:strong:electrovac}

A strong gravitational field is a useful test of a constrained construction because coordinate singularities, physical curvature, and numerical constraint drift can occur at quite different places. The spherical models below keep those distinctions visible. The electrovac model supplies an exact curved spacetime whose coordinates move in time. The scalar model supplies dynamical matter and a numerically observed marginally trapped sphere. Their computational maps have different locality properties, and their evidence supports different conclusions.

The electrovac sector sets the Dirac field to zero and imposes spherical symmetry. It retains a real canonical Einstein--Maxwell system, but it has no nonspherical gravitational radiation. Work in units $G=c=4\pi\epsilon_0=1$. A prime denotes a radial derivative and a dot denotes a time derivative. On a radial interval where $R>0$ and $\Lambda>0$, write
\begin{equation}
 ds^2=-N^2dt^2+\Lambda^2(dr+N^r dt)^2+R^2d\Omega^2,
 \qquad \mathcal A=A_rdr+\Phi_e dt.
\end{equation}
Here $R$ is the areal radius, $\Lambda$ determines proper radial distance, $N$ is the lapse, and $N^r$ is the radial shift. The Maxwell multiplier is denoted $\Phi_e$ to distinguish it from the scalar gradient used later. The bulk canonical action is
\begin{equation}
 S=\int dt\int dr\,
 \bigl(P_\Lambda\dot\Lambda+P_R\dot R+Q\dot A_r
       -N\mathcal H-N^r\mathcal D-\Phi_e\mathcal G\bigr)+B_\partial.
\end{equation}
The boundary term specifies the mass and electromagnetic boundary ensemble. Its treatment cannot be omitted when claiming a global Hamiltonian evolution, even though it does not change the following local bulk identity.

The gauge-improved radial momentum and electric constraints are
\begin{equation}
 \mathcal D=P_RR'-\Lambda P_\Lambda',\qquad \mathcal G=Q',
\end{equation}
and the Hamiltonian constraint is
\begin{equation}
\begin{aligned}
 \mathcal H={}&-\frac{P_RP_\Lambda}{R}
 +\frac{\Lambda P_\Lambda^2}{2R^2}
 +\frac{RR''}{\Lambda}-\frac{RR'\Lambda'}{\Lambda^2}
 +\frac{R'^2}{2\Lambda}-\frac\Lambda2
 +\frac{\Lambda Q^2}{2R^2}.
\end{aligned}
\label{c:strong:electrovac-H}
\end{equation}
The momenta follow the convention
\begin{align}
 P_\Lambda&=-\frac RN(\dot R-N^rR'),\\
 P_R&=-\frac1N\left[\Lambda(\dot R-N^rR')
                  +R\bigl(\dot\Lambda-(N^r\Lambda)'\bigr)\right].
\end{align}
These signs fix which Painlev\'e--Gullstrand branch describes inward future characteristics.

\begin{definition}[Charged mass coordinate]
The charged Misner--Sharp coordinate in this canonical chart is
\begin{equation}
 \mathcal M=\frac R2\left[1-\left(\frac{R'}\Lambda\right)^2
                    +\left(\frac{P_\Lambda}R\right)^2\right]
             +\frac{Q^2}{2R}.
\label{c:strong:mass-coordinate}
\end{equation}
The charge contribution is included in $\mathcal M$ so that it is constant throughout an electrovac Reissner--Nordstr\"om solution.
\end{definition}

\begin{proposition}[Exact triangular mass identity]
For differentiable canonical data in the regular chart,
\begin{equation}
 \mathcal M'=-\frac{R'}\Lambda\mathcal H
             -\frac{P_\Lambda}{R\Lambda}\mathcal D
             +\frac QR\mathcal G.
\label{c:strong:mass-identity}
\end{equation}
Consequently, on a patch where $R'/\Lambda\ne0$, the systems
$(\mathcal H,\mathcal D,\mathcal G)=0$ and
$(\mathcal M',\mathcal D,\mathcal G)=0$ have the same zero set.
\end{proposition}
\begin{proof}
Differentiating \eqref{c:strong:mass-coordinate} gives
\begin{equation}
\begin{aligned}
 \mathcal M'={}&\frac{R'}2-\frac{R'^3}{2\Lambda^2}
 -\frac{RR'R''}{\Lambda^2}+\frac{RR'^2\Lambda'}{\Lambda^3}
 +\frac{P_\Lambda P_\Lambda'}R-\frac{P_\Lambda^2R'}{2R^2}\\
 &+\frac{QQ'}R-\frac{Q^2R'}{2R^2}.
\end{aligned}
\end{equation}
Expanding the right side of \eqref{c:strong:mass-identity} cancels the two $P_R$ terms and produces precisely this expression. The constraint transformation is upper triangular with diagonal $(-R'/\Lambda,1,1)$, so it is invertible under the stated condition.
\end{proof}

This calculation changes constraint coordinates before discretization. It does not prove closure of an arbitrary pointwise discretization of the original hypersurface-deformation algebra. It also establishes equivalence of zero sets, rather than proving that every Poisson bracket of the replacement functions vanishes. The distinction matters when the reduced construction is compared with a general BFV system.

\section{Horizon penetration and the regular radial chart}
\label{c:strong:RN-chart}

Fix $M>0$, a real charge $Q$, and a branch $\varepsilon\in\{+1,-1\}$. On $2MR-Q^2>0$, define
\begin{equation}
 p(R)=\varepsilon\sqrt{2MR-Q^2},\qquad v(R)=\frac{p(R)}R.
\end{equation}
The Painlev\'e--Gullstrand metric takes the form
\begin{equation}
 ds^2=-dt_P^2+(dR+v(R)dt_P)^2+R^2d\Omega^2.
\label{c:strong:RN-PG}
\end{equation}
No inverse factor of $1-2M/R+Q^2/R^2$ appears. The metric can therefore cross a horizon while the lapse and radial spatial metric remain regular.

Take a smooth radial map $R=F(t,r)$ with $F_r>0$ and identify $t_P=t$. Expanding $dR=F_tdt+F_rdr$ identifies
\begin{equation}
 N=1,\qquad \Lambda=F_r,\qquad N^r=\frac{F_t+v(F)}{F_r}.
\end{equation}
The corresponding canonical data are
\begin{equation}
 P_\Lambda=p(F),\qquad P_R=\frac{MF_r}{p(F)},
 \qquad A_r=0,\quad \Phi_e=-\frac QF,\quad Q'=0.
\label{c:strong:RN-data}
\end{equation}
The sign in $P_R$ is contained in $p(F)$.

\begin{proposition}[A time-dependent exact gauge family]
The data \eqref{c:strong:RN-data} satisfy the local electrovac constraints and canonical evolution equations wherever $F_r>0$, $F>0$, and $2MF-Q^2>0$.
\end{proposition}
\begin{proof}
First, $R'/\Lambda=1$ and $p(F)^2=2MF-Q^2$ imply $\mathcal M=M$. Since
\begin{equation}
 P_\Lambda'=\frac{MF_r}{p(F)}=P_R,
\end{equation}
the momentum constraint vanishes. The mass identity and $Q'=0$ then give $\mathcal H=0$.

For evolution, the radius equation gives
\begin{equation}
 \dot R=-\frac{P_\Lambda}R+N^rR'
        =-v(F)+F_t+v(F)=F_t.
\end{equation}
The radial-metric equation is
\begin{equation}
 \dot\Lambda=\frac{\Lambda P_\Lambda}{R^2}
             -\frac{P_R}R+(N^r\Lambda)'.
\end{equation}
Use $p_R=M/p$ and $v_R=M/(pR)-p/R^2$. The non-$F_{tr}$ terms cancel, leaving $\dot\Lambda=F_{tr}$. The momentum equations in the $N=1$ convention are
\begin{align}
 \dot P_\Lambda&=N^rP_\Lambda'
 +\frac12\left[1-\frac{P_\Lambda^2+Q^2}{R^2}
                 -\left(\frac{R'}\Lambda\right)^2\right],\\
 \dot P_R&=\frac{\Lambda(P_\Lambda^2+Q^2)}{R^3}
 -\frac{P_\Lambda P_R}{R^2}+(N^rP_R)'
 -\frac{R''}\Lambda+\frac{\Lambda'R'}{\Lambda^2}.
\end{align}
Substituting the data and differentiating $p(F)^2=2MF-Q^2$ reduces their right sides to
\begin{equation}
 \frac{MF_t}{p(F)},\qquad
 M\left[\frac{F_{tr}}{p(F)}
       -\frac{MF_rF_t}{p(F)^3}\right],
\end{equation}
which are the direct time derivatives of the two momenta. Finally,
$\Phi_e'=QF_r/F^2$ cancels $N\Lambda Q/R^2$ in
$\dot A_r=N\Lambda Q/R^2-\Phi_e'$.
\end{proof}

An arbitrary admissible $F$ therefore creates coordinate motion through one fixed electrovac spacetime. Neither the mass nor the charge changes. A numerical reconstruction of this family is a nonlinear strong-field consistency test, but it is not matter collapse.

\subsection{Charged horizons and the turning surface}
Write
\begin{equation}
 f(R)=1-\frac{2M}R+\frac{Q^2}{R^2}.
\end{equation}
For $0<|Q|<M$, its two positive roots are
\begin{equation}
 r_\pm=M\pm\sqrt{M^2-Q^2}.
\end{equation}
At either root $p(R)^2=R^2$, so the canonical quantities in \eqref{c:strong:RN-data} remain finite if $F_r$ does. For $|Q|=M$, the two roots coincide at $R=M$. For $|Q|>M$ there is no real root. At $Q=0$, the outer root is $2M$ and the formal inner root $0$ is the singular endpoint, not a second regular horizon.

The square-root chart instead meets a turning surface at
\begin{equation}
 R_{\rm turn}=\frac{Q^2}{2M}.
\end{equation}
There $p=0$, and $P_R=MF_r/p$ becomes unbounded for nonzero $F_r$. This is a limitation of the chosen flat spatial slice and canonical chart. Horizon penetration does not remove the need for a positive distance from this turning surface. It also does not remove the physical curvature singularity at $R=0$.

The invariant curvature diagnostic is
\begin{equation}
 \mathcal K=\frac{48M^2}{R^6}-\frac{96MQ^2}{R^7}
                            +\frac{56Q^4}{R^8}.
\end{equation}
For $M=1$, $Q=0.8$, the radii are $r_+=1.6$, $r_-=0.4$, and $R_{\rm turn}=0.32$. The benchmark interval $R\in[0.6,6]$ crosses the outer horizon only. At its inner end $\mathcal K\simeq199.6646853$, while $\mathcal K(r_+)\simeq1.106262207$. High curvature and regular horizon-crossing coordinates coexist on this interval.

\gravityfigure{10_electrovac_horizons}{The horizon equation and the domain of the charged Painlev\'e--Gullstrand chart are different conditions. The benchmark crosses the outer horizon while remaining above the turning surface. The extremal and superextremal curves illustrate the change in the number and multiplicity of horizon roots.}

\subsection{Horizon conditioning}
For a subextremal outer horizon, put $d=\sqrt{M^2-Q^2}>0$. Then
\begin{equation}
 f'(r_+)=\frac{2d}{r_+^2}.
\end{equation}
If a nearby continuous numerical profile has a selected root near this simple root, the local inverse estimate gives
\begin{equation}
 |\delta r_+|\le\frac{r_+^2}{2d}\,\norm{\delta f}_\infty
                    +O(\norm{\delta f}_\infty^2).
\label{c:strong:root-bound}
\end{equation}
The estimate describes the nearby root and presupposes its localization to a region where the exact derivative is nonzero. It does not make every zero of an arbitrarily oscillatory perturbation physically relevant. The factor is $2.133333333\ldots$ in the benchmark and diverges as $d\downarrow0$.

At extremality $f(R)=(1-M/R)^2$ has a double root. A perturbation of size $\epsilon$ can move the root by order $\sqrt\epsilon$ or remove it entirely. Finite canonical data at a horizon therefore do not imply uniformly well-conditioned horizon detection. The first trapped sphere in the scalar calculation below also occurs at a spatial compactness maximum, and must not be assigned this simple-root estimate without checking its event geometry.

\section{Finite radial reconstruction and its stability boundary}
\label{c:strong:finite-reconstruction}

On interior points of a uniform radial grid, use the fourth-order derivative
\begin{equation}
 (D_4u)_j=\frac{-u_{j+2}+8u_{j+1}-8u_{j-1}+u_{j-2}}{12h}.
\end{equation}
Its stencil radius is two. Given $M$, $Q$, and a fixed branch $\varepsilon$, a radius array satisfying $D_4R_j>0$ and $2MR_j-Q^2>0$ determines
\begin{equation}
 \Pi_h[R]_j=\left(R_j,D_4R_j,
                 \varepsilon\sqrt{2MR_j-Q^2},
                 \frac{M D_4R_j}{\varepsilon\sqrt{2MR_j-Q^2}},Q\right).
\label{c:strong:reconstruct}
\end{equation}
This is a local reconstruction of $(R,\Lambda,P_\Lambda,P_R,Q)$, rather than an orthogonal Hilbert-space projector. If it is extended to a full state by discarding the other variables and retaining $R$, it is idempotent on its admissible image.

\Needspace{18\baselineskip}
\begin{proposition}[Exact finite reduced constraints]
On the image of \eqref{c:strong:reconstruct}, the four reduced constraints
\begin{equation}
\begin{aligned}
 C_{\Lambda,j}&=\Lambda_j-D_4R_j,&
 C_{M,j}&=P_{\Lambda,j}^2-2MR_j+Q^2,\\
 C_{P,j}&=P_{R,j}P_{\Lambda,j}-M\Lambda_j,&
 C_{Q,j}&=Q_j-Q
\end{aligned}
\end{equation}
vanish in exact arithmetic. The discrete mass coordinate obtained by replacing $R'$ with $D_4R$ in \eqref{c:strong:mass-coordinate} equals $M$ at every reconstructed node.
\end{proposition}
\begin{proof}
Each constraint vanishes by direct substitution. In the discrete mass coordinate $D_4R/\Lambda=1$ and $P_\Lambda^2=2MR-Q^2$, leaving
$\tfrac{R}{2}(2M/R-Q^2/R^2)+Q^2/(2R)=M$.
\end{proof}

These identities are exact algebraic constraints on the reconstruction manifold. The original ADM expressions evaluated with finite differences have a different status: the stencil does not obey a pointwise chain rule. Thus $D_4\sqrt{2MR-Q^2}$ need not equal $MD_4R/\sqrt{2MR-Q^2}$ exactly. The original momentum and Hamiltonian constraints are retained as continuum-consistency diagnostics, rather than being relabeled as exact finite invariants.

\subsection{A level-independent reconstruction estimate}
Assume a uniform regular set
\begin{equation}
 R_j\ge R_*>0,\quad 2MR_j-Q^2\ge p_*^2>0,
 \quad 0<d_*\le D_4R_j\le d^*.
\end{equation}
For two arrays in this set, with the same $(M,Q,\varepsilon)$, differentiation along the segment between them gives
\begin{align}
 \norm{\delta\Lambda}_\infty&\le\norm{D_4\delta R}_\infty,\\
 \norm{\delta P_\Lambda}_\infty&\le\frac M{p_*}\norm{\delta R}_\infty,\\
 \norm{\delta P_R}_\infty&\le\frac M{p_*}\norm{D_4\delta R}_\infty
                          +\frac{M^2d^*}{p_*^3}\norm{\delta R}_\infty.
\label{c:strong:projection-bound}
\end{align}
For example, $\partial_Rp=M/p$ and the two partial derivatives of $M\Lambda/p$ have magnitudes at most $M/p_*$ and $M^2d^*/p_*^3$. The admissible segment preserves the displayed linear radius and derivative bounds. The constants are independent of $h$.

The input norm contains a radial derivative. A small pointwise radius error alone therefore does not establish a small error in every reconstructed canonical field. This estimate transfers an already controlled discrete $W^{1,\infty}$ radius error to the reconstructed state. It does not create that derivative control from an $L^2$ estimate.

\subsection{Direct and split radius evolution}
On the positive branch, restriction of the radius equation to the reconstruction manifold gives
\begin{equation}
 \dot R_j=g(R_j)+b_j(D_4R)_j,
 \qquad g(R)=-\frac{\sqrt{2MR-Q^2}}R,
\label{c:strong:direct-radius}
\end{equation}
where $b=N^r$ is prescribed in the electrovac benchmark. A split alternative is
\begin{equation}
 \dot R=g(R)+A_h[b]R,\qquad
 A_h[b]u=\frac12\bigl(bD_4u+D_4(bu)-(D_4b)u\bigr).
\label{c:strong:split-radius}
\end{equation}
On smooth data these formulas differ by half of the discrete product-rule defect, which is fourth order away from boundaries. Their exact finite-grid operators differ.

The archived implementation advances the radius explicitly with RK4. The map $\Pi_h$ can reconstruct the other fields from each admissible stage radius, and its reduced identities hold on that reconstructed image. At output, however, the code overwrites two canonical boundary nodes on each side with exact manufactured data. Those boundary nodes are not generally the exact image of the finite derivative reconstruction. Its reported constraint norms exclude four nodes from each end. This procedure does not use an implicit discrete-gradient slab solve. Conservation by an implicit invariant-preserving method elsewhere in this monograph remains conditional on actually solving its nonlinear step equation. It cannot be used as evidence that this explicit radius algorithm preserves the original ADM Hamiltonian exactly.

\begin{proposition}[Split-form energy identity]
Let $D$ be skew adjoint for the grid inner product $\inner{u}{w}_h$, and let multiplication by the real array $b$ be self-adjoint. Then
\begin{equation}
 \inner{u}{A_h[b]u}_h
       =-\frac12\inner{u}{(Db)u}_h.
\label{c:strong:split-energy}
\end{equation}
For two solutions of the split equation with the same prescribed $b$, remaining in a positive-radicand interval where $|g'|\le L_g$, one obtains
\begin{equation}
 \norm{R(t)-\widetilde R(t)}_h
 \le e^{(L_g+S_1/2)t}\norm{R(0)-\widetilde R(0)}_h,
 \qquad \norm{Db}_\infty\le S_1.
\label{c:strong:semidiscrete-bound}
\end{equation}
\end{proposition}
\begin{proof}
Skew adjointness cancels $\inner{u}{bDu}_h+\inner{u}{D(bu)}_h$. The remaining term gives \eqref{c:strong:split-energy}. With $e=R-\widetilde R$, the nonlinear difference satisfies
$|g(R_j)-g(\widetilde R_j)|\le L_g|e_j|$, where
\begin{equation}
 g'(R)=\frac{MR-Q^2}{R^2\sqrt{2MR-Q^2}}.
\end{equation}
Thus $\tfrac12\partial_t\norm e_h^2\le(L_g+S_1/2)\norm e_h^2$, and integration gives the claim.
\end{proof}

The skew-adjoint hypothesis needs care on a physical radial interval. A periodic $D_4$ satisfies it, but a globally periodic radius array cannot also satisfy $D_4R_j>0$ everywhere: summing the derivative over one period gives zero. The periodic identity is an exact algebraic check of the split stencil. It does not by itself give a nonvacuous stability theorem for the globally monotone spherical reconstruction.

For an interval summation-by-parts derivative with a diagonal positive norm matrix $H$ and $HD+D^TH=B$, multiplication by $b$ commutes with $H$. The uncancelled term is instead
\begin{equation}
 \inner{u}{A_h[b]u}_H
 =\frac12u^TB\,\operatorname{diag}(b)u
  -\frac12\inner{u}{(Db)u}_H.
\end{equation}
Its boundary flux must be controlled using an actual boundary closure and appropriate boundary data. Saying that a stencil is centered in the interior does not establish this property for the complete interval operator.

The periodic Fourier symbol is
\begin{equation}
 \widehat D_4(\theta)=\frac{i}{h}\frac{8\sin\theta-\sin2\theta}{6}.
\end{equation}
Its maximum dimensionless magnitude is $1.3722219798\ldots$. Applying the RK4 imaginary-axis stability interval to frozen constant advection gives
\begin{equation}
 \frac{|b|\Delta t}{h}\le2.0612023174\ldots.
\end{equation}
This is a frozen-coefficient diagnostic. A full variable-coefficient nonlinear RK4 convergence theorem still requires a uniform stability argument for the implemented boundary closure and a suitable stronger norm. Fourth-order consistency, the semidiscrete $L^2$ estimate, and this one CFL number do not alone prove fourth-order $W^{1,\infty}$ convergence.

\section{What the electrovac benchmark measures}
\label{c:strong:electrovac-evidence}

The retained run uses
\begin{equation}
 F(t,r)=r+0.08\sin(0.7t)
                \sin^2\!\left(\pi\frac{r-0.6}{5.4}\right),
 \quad 0\le t\le0.4,
\end{equation}
with $M=1$, $Q=0.8$, and $\Delta t/h$ close to $0.1$. The sampled minimum radial Jacobian is $0.9871380683$ and the minimum momentum radicand is $0.56$. The exact shift, ghost values, and two boundary nodes on either side are supplied from the manufactured family. Errors are measured on interior nodes after excluding four points from each end.

The recorded state errors are
\begin{center}\small
\begin{tabular}{rrrrr}
\toprule
$N$&$\norm{\delta R}_\infty$&$\norm{\delta\Lambda}_\infty$&$\norm{\delta P_\Lambda}_\infty$&$\norm{\delta P_R}_\infty$\\
\midrule
64&$8.3247\,10^{-9}$&$4.1362\,10^{-8}$&$5.7367\,10^{-9}$&$2.3732\,10^{-8}$\\
128&$5.0495\,10^{-10}$&$2.5079\,10^{-9}$&$3.3079\,10^{-10}$&$1.3795\,10^{-9}$\\
256&$3.1077\,10^{-11}$&$1.5438\,10^{-10}$&$2.0248\,10^{-11}$&$8.4854\,10^{-11}$\\
512&$1.9296\,10^{-12}$&$1.0485\,10^{-11}$&$1.2559\,10^{-12}$&$5.3844\,10^{-12}$\\
\bottomrule
\end{tabular}
\end{center}
The least-squares orders are respectively $3.99856$, $3.96010$, $4.02393$, and $4.00788$. The ordinary ADM Hamiltonian on exact projected data has fitted order $3.70410$. These are results for a smooth, forced trajectory on four grids. The tiny final errors are close enough to roundoff that they should not be interpreted as arbitrarily extendible asymptotic scaling laws.

On the tested interior of the $N=512$ reconstructed grid, the retained mass residual is $4.44\,10^{-16}$, the flat-slice residual is zero, the algebraic mass constraint is $2.22\,10^{-15}$, and the reduced momentum residual is $1.11\,10^{-16}$. These floating-point numbers are consistent with the exact algebraic reconstruction theorem. They are not an exact-zero assertion about floating arithmetic or the overwritten boundary nodes.

The negative control evolves all four canonical fields without the reconstruction. At $N=128$, it reports mass drift $2.25125\,10^{-4}$, radial momentum residual $1.27688\,10^{-3}$, and Hamiltonian residual $1.25735\,10^{-2}$. This gives the reduced reconstruction a meaningful comparator. It does not show that every other constrained algorithm fails or that the direct ADM system is ill posed.

A smooth perturbation of amplitude $10^{-8}$ evolved with the split form has recorded $L^2$ amplification between $1.03345462$ and $1.03348106$ on $N=128,256,512,1024$. The spread is $2.64438\,10^{-5}$. This supports the bounded experiment's level-insensitive behavior. It leaves the interval boundary and stronger-norm theorem obligations visible.

The outer horizon remains at the fixed areal radius $1.6$. Its computational coordinate moves because $F$ changes. The retained final coordinate comparison at $N=256$ gives $|r_{+,h}-r_+^{\rm exact}|=1.89262\,10^{-11}$. This measures tracking of a known horizon under coordinate deformation. It does not measure the formation or growth of a black hole.

The electrovac verification record contains nineteen passed checks. The symbolic identity and archived data structure have supporting audit records, but these are distinct from a replay of the electrovac evolution. The numerical table is retained as recorded evidence. Its environment and command provenance are specified in the verification appendix.

\section{A dynamical spherical scalar system}
\label{c:strong:scalar-reduction}

To test dynamical concentration of matter, replace electrovac by a real massless scalar. Use $G=c=1$ and the rescaled field $\varphi=\sqrt{4\pi}\,\phi_{\rm phys}$. The scalar canonical pair is $(\varphi,P_\varphi)$, and the matter additions to the Hamiltonian and momentum constraints are
\begin{equation}
 \mathcal H_{\rm matt}=\frac{P_\varphi^2}{2\Lambda R^2}
                         +\frac{R^2\varphi'^2}{2\Lambda},
 \qquad \mathcal D_{\rm matt}=P_\varphi\varphi'.
\end{equation}
The gravitational part is \eqref{c:strong:electrovac-H} with $Q=0$.

Impose $R=r$ and $\Lambda=1$, and set
\begin{equation}
 N=\sigma>0,\quad N^r=\sigma v,\quad
 P_\Lambda=rv,\quad P_\varphi=r^2\Pi,\quad \Phi=\varphi'.
\end{equation}
Here $\Pi$ is the scalar momentum per areal area, not the reconstruction map $\Pi_h$. Define
\begin{equation}
 m=\frac r2v^2,\qquad v=\sqrt{\frac{2m}{r}}.
\end{equation}
The inward branch gives
\begin{equation}
 ds^2=-\sigma^2dt^2+(dr+\sigma vdt)^2+r^2d\Omega^2.
\label{c:strong:scalar-metric}
\end{equation}
The following derivation applies on smooth patches with $r>0$, $\sigma>0$, and $v>0$. These hypotheses keep every division legitimate.

The momentum constraint removes $P_R$ algebraically:
\begin{equation}
 P_R=P_\Lambda'-P_\varphi\Phi.
\end{equation}
Since $m'=P_\Lambda P_\Lambda'/r-P_\Lambda^2/(2r^2)$, inserting this relation into the Hamiltonian constraint gives
\begin{equation}
 m'=r^2\mathcal E,\qquad
 \mathcal E=\frac12(\Pi^2+\Phi^2)+v\Pi\Phi.
\label{c:strong:scalar-mass}
\end{equation}
Preservation of $\Lambda=1$ gives
\begin{equation}
 0=\sigma\left(\frac vr-\frac{P_R}r\right)+(\sigma v)'
   =\sigma r\Pi\Phi+\sigma'v,
\end{equation}
hence
\begin{equation}
 (\log\sigma)'=-\frac{r\Pi\Phi}{v}.
\label{c:strong:scalar-lapse}
\end{equation}
The scalar evolution equations are
\begin{align}
 \dot\varphi&=\sigma(\Pi+v\Phi),\\
 \dot\Phi&=\partial_r[\sigma(\Pi+v\Phi)],\\
 \dot\Pi&=r^{-2}\partial_r[r^2\sigma(\Phi+v\Pi)].
\label{c:strong:scalar-evolution}
\end{align}
If $\Phi=\varphi'$ initially and compatible boundary data are imposed, the first two equations preserve that relation in the continuum.

At $v=0$, the lapse equation requires a separately specified extension. Smoothness of $\Pi$ and $\Phi$ alone does not establish that $r\Pi\Phi/v$ is removable. A vacuum interval with vanishing scalar data can provide an appropriate extension, but that must be part of the argument. The Gaussian benchmark has very small tails rather than exact compact support, and its numerical treatment of small $v$ is described below. This calculation does not prove a regular-center theorem.

\subsection{Exact mass flux and compatibility}
Put $A=\Pi+v\Phi$ and $B=\Phi+v\Pi$. The remaining Einstein momentum evolution gives
\begin{equation}
 \dot m=r^2\mathcal J,\qquad \mathcal J=\sigma AB.
\label{c:strong:scalar-flux}
\end{equation}
One can verify its origin directly from the $P_\Lambda$ equation. In the flat-slice gauge, the gravitational terms give $\sigma vP_\Lambda'-\sigma v^2/2-r\sigma'$. The lapse-gradient term is essential. The matter variation contributes $\sigma r^2(\Pi^2+\Phi^2)/2$. As $\dot m=v\dot P_\Lambda$, their sum becomes
\begin{equation}
 \dot m=\sigma v\left[m'+\frac{r^2}{2}(\Pi^2+\Phi^2)\right]-rv\sigma'.
\end{equation}
Using \eqref{c:strong:scalar-mass} and \eqref{c:strong:scalar-lapse} yields
$\dot m=r^2\sigma[v(\Pi^2+\Phi^2)+(1+v^2)\Pi\Phi]$, which is \eqref{c:strong:scalar-flux}.

\begin{proposition}[Continuum mass-constraint propagation]
On the regular patch, the scalar evolution, lapse relation, and mass-flux evolution satisfy
\begin{equation}
 \partial_t(r^2\mathcal E)=\partial_r(r^2\mathcal J).
\label{c:strong:scalar-compatibility}
\end{equation}
Thus $\partial_t(m'-r^2\mathcal E)=0$. A radial mass constraint imposed initially propagates under this continuum system.
\end{proposition}
\Needspace{20\baselineskip}
\begin{proof}
Expanding \eqref{c:strong:scalar-evolution} gives
$\dot\Phi=(\sigma A)'$ and $\dot\Pi=(\sigma B)'+2\sigma B/r$.
Differentiate $\mathcal E$ and subtract the radial derivative of the flux. Terms involving $\Pi'$, $\Phi'$, and $v'$ cancel, leaving
\begin{equation}
 \partial_t(r^2\mathcal E)-\partial_r(r^2\mathcal J)
 =r^2\bigl[\sigma'AB+\Pi\Phi\dot v\bigr].
\end{equation}
Since $m=rv^2/2$, the flux equation implies $\dot v=r\sigma AB/v$. The remaining expression is
\begin{equation}
 r^2AB\left[\sigma'+\frac{\sigma r\Pi\Phi}{v}\right]=0
\end{equation}
by \eqref{c:strong:scalar-lapse}. This establishes the identity without assuming the radial constraint in its last cancellation step.
\end{proof}

On a fixed radial interval $[r_i,r_o]$, the exact balance is
\begin{equation}
 \frac{d}{dt}\bigl(m(t,r_o)-m(t,r_i)\bigr)
 =r_o^2\mathcal J(t,r_o)-r_i^2\mathcal J(t,r_i).
\end{equation}
Boundary fluxes, rather than a statement that a finite-domain mass must be exactly constant, are the appropriate continuum reference. The numerical prefix solve below does not satisfy this identity algebraically at every finite grid and time step.

\subsection{Characteristic channels and the loss of a coordinate energy bound}
Define $w_+=\Pi+\Phi$ and $w_-=\Pi-\Phi$. Their equations are
\begin{align}
 \dot w_+={}&\sigma(v+1)w_+'
 +[\sigma v'+\sigma'(v+1)]w_+
 +\frac\sigma r[(v+1)w_++(v-1)w_-],\\
 \dot w_-={}&\sigma(v-1)w_-'
 +[\sigma v'+\sigma'(v-1)]w_-
 +\frac\sigma r[(v-1)w_-+(v+1)w_+].
\end{align}
Written as transport equations $\dot w+c\,w'=\text{source}$, their coordinate speeds are
\begin{equation}
 c_+=-\sigma(v+1),\qquad c_-=\sigma(1-v).
\label{c:strong:characteristics}
\end{equation}
The plus channel is inward on the positive branch. The minus channel is outward when $v<1$, stationary at $v=1$, and inward when $v>1$.

The same channels give the exact algebraic identities
\begin{align}
 \mathcal E&=\frac14[(1+v)w_+^2+(1-v)w_-^2],\\
 \mathcal J&=\frac\sigma4[(1+v)^2w_+^2-(1-v)^2w_-^2].
\label{c:strong:null-energy}
\end{align}
For $0\le v<1$, $\mathcal E\ge(1-v)(w_+^2+w_-^2)/4$. This lower bound degenerates at a horizon. For $v>1$, the outward-channel coefficient in this coordinate mass density is negative. The scalar's local normal-frame quadratic density $(\Pi^2+\Phi^2)/2$ remains nonnegative. The loss of coercivity in $\mathcal E$ therefore diagnoses this slicing's mass-coordinate energy estimate, rather than negative physical scalar energy.

\gravityfigure{11_scalar_flux}{The mass density and mass flux resolve into the two scalar characteristic channels. The outgoing characteristic and its contribution to the coordinate mass density degenerate at $v=1$. These are exact algebraic profiles, not a newly evolved collapse solution.}

\section{Trapped spheres and the finite-stencil scalar experiment}
\label{c:strong:collapse-evidence}

For \eqref{c:strong:scalar-metric}, choose future radial null vectors
\begin{equation}
 \ell_\pm=\sigma^{-1}\partial_t+(-v\pm1)\partial_r.
\end{equation}
Because the sphere area is $4\pi r^2$, its expansion along $\ell$ is $2\ell(r)/r$. Consequently
\begin{equation}
 \theta_+=\frac2r(1-v),\qquad
 \theta_-=-\frac2r(1+v).
\label{c:strong:expansions}
\end{equation}
On the positive branch, $v=1$ is a marginally outer trapped sphere, and $v>1$ gives both expansions negative. Equivalently, the compactness $2m/r$ equals or exceeds one. The normalization of the null vectors affects expansion magnitude but not these signs. A marginally trapped sphere is a local geometric diagnostic. It is not by itself identification of a global event horizon.

\subsection{What is local in the algorithm}
At each RK4 stage the matter channels determine the geometry through two radial solves. The mass constraint is integrated outward with $m(r_{\rm exc})=0$. The lapse derivative is integrated and the result normalized by $\sigma(r_{\max})=1$. These are triangular prefix operations. A matter value near the inner end can affect the reconstructed mass and lapse far along the interval in the same stage. The full stage map is therefore globally coupled.

The characteristic evolution uses the second-order one-sided derivatives
\begin{equation}
 (D_+u)_j=\frac{-3u_j+4u_{j+1}-u_{j+2}}{2h},\qquad
 (D_-u)_j=\frac{3u_j-4u_{j-1}+u_{j-2}}{2h}.
\end{equation}
The inward channel uses $D_+$ and the outgoing channel uses $D_-$ outside a horizon, switching to $D_+$ inside. These interior matter stencils have radius two. The source also uses a radius-two fourth-order interior derivative for $v'$ and separate one-sided boundary formulas. Finite stencil radius for these derivatives must not be promoted to a finite dependency cone for the geometry-coupled step.

The mass integration uses radial RK4 with interpolated midpoint fields, followed by Simpson-type lapse quadrature. These are approximate quadratures. They are neither exact mass-constraint algebraic reconstruction nor an implicit invariant-preserving solver. The scalar spatial evolution is second order despite its fourth-order time integrator.

Several implementation conventions are mathematically material. The mass right side and the reconstructed $v$ use $\sqrt{\max(2m/r,0)}$. The lapse derivative is set to zero where $v\le10^{-13}$. Incoming characteristic values are explicitly set after a completed RK4 step, with separate boundary formulas during stages. These choices guard a finite calculation near vacuum and boundaries. An exact continuum proof on $v>0$ does not establish their uniform consistency as the threshold is varied or $h\downarrow0$.

\subsection{Initial data and boundary scope}
The numerical domain is $[r_{\rm exc},r_{\max}]=[0.05,30]$. It replaces a regular center by a small excised core. The characteristic data are $w_-(r_{\rm exc})=0$ and $w_+(r_{\max})=0$. These prohibit prescribed incoming radiation from the inner and outer boundaries in the respective channels. The radial mass boundary condition additionally fixes the excised mass to zero at every geometry solve.

The initial scalar packet is
\begin{equation}
 \varphi(0,r)=A e^{-(r-10)^2},\qquad
 \Phi(0,r)=\partial_r\varphi(0,r),\qquad
 \Pi(0,r)=\Phi(0,r)+\frac{\varphi(0,r)}r.
\end{equation}
The last equation is the flat-space ingoing condition for $\chi=r\varphi$, used to initialize the scalar. The nonlinear gravitational constraints are then solved. It is therefore an ingoing packet prescription, rather than exactly freezing the geometry to Minkowski space.

The strong amplitude is $A=0.10$ and the control amplitude is $A=0.05$. The strong event occurs near $r=3.1$, well separated from both boundaries in coordinate distance. That separation is useful but is not a substitute for studying boundary influence: the geometry solve is global, and finite-domain mass drift can contain both boundary and discretization effects.

\subsection{Recorded first marginally trapped sphere}
The event finder takes a quadratic interpolation of the spatial maximum of $2m/r$ and a linear interpolation between consecutive time steps bracketing its first crossing of one. The archived event table is
\begin{center}\small
\begin{tabular}{rrrrr}
\toprule
$N_r$&$h$&$t_H$&$r_H$&$\Delta M/M_0$\\
\midrule
600&0.0500000&4.7133579&3.0518959&$-1.1205\,10^{-2}$\\
800&0.0374844&4.7162491&3.0506543&$-1.6427\,10^{-2}$\\
1000&0.0299800&4.7059219&3.0766099&$-1.0696\,10^{-2}$\\
1200&0.0249791&4.7004729&3.0905271&$-7.6261\,10^{-3}$\\
1600&0.0187305&4.6952990&3.1040860&$-4.6604\,10^{-3}$\\
2000&0.0149825&4.6930194&3.1101678&$-3.3322\,10^{-3}$\\
2400&0.0124844&4.6918209&3.1134115&$-2.6241\,10^{-3}$\\
\bottomrule
\end{tabular}
\end{center}
The coarsest two event coordinates are not monotone with refinement. On the six finest resolutions the retained nonlinear fits are
\begin{align}
 t_H(h)&\simeq4.6893522370+33.88577338\,h^{2.173853282},\\
 r_H(h)&\simeq3.1204965744-65.59310603\,h^{2.084394965}.
\end{align}
These fit exponents are consistent with second-order event convergence on this bounded family. They are not a convergence proof, a reliable error bar for every observable, or a critical-collapse exponent.

For $N_r=2400$, the interpolated event has
\begin{equation}
 t_H=4.69182094580,\quad r_H=3.11341154797,
 \quad m_H=r_H/2=1.55670577399.
\end{equation}
The relation $m_H=r_H/2$ is imposed by the interpolated marginal-sphere criterion. It is not an independent numerical conservation test.

The stored event fields come from the next completed step, at $t=4.69206969694$, rather than exactly the interpolated $t_H$. Their recorded outer mass is $1.72538868709$, compared with initial outer mass $1.72992821318$, a fractional change of $-2.6241124\,10^{-3}$. The same stored fields have minimum lapse $0.999829382946$, lapse at the sampled peak $1.02582462096$, and maximum absolute scalar Ricci curvature $7.35797274672$. This timing distinction is small but matters when those quantities are compared with an independently interpolated event.

The scalar curvature diagnostic follows the rescaled convention
\begin{equation}
 \mathcal R=2(\Phi^2-\Pi^2).
\end{equation}
Finite lapse and curvature values show that this calculation can reach the marginal-sphere event without a coordinate blow-up. They do not provide continuation to a singularity or control of subsequent trapped-region evolution.

\gravityfigure{12_trapped_refinement}{Archived event coordinates under radial refinement, together with the exact expansion and characteristic sign change at $v=1$. The event table is replotted from retained data. No new scalar-collapse evolution was performed for this figure.}

\subsection{Control and independent diagnostics}
The weak pulse on $N_r=1200$ remains untrapped through $t=5$, with maximum compactness $0.146722872367$ and fractional outer-mass change $-4.5106042\,10^{-3}$. It is a finite-time control in an excised domain. It does not establish global dispersion, passage through a regular center, or a critical amplitude separating two complete physical outcomes.

At the smooth pre-horizon time $t=3$, the $N_r=1200$ retained relative residuals are
\begin{center}
\begin{tabular}{lr}
\toprule
Diagnostic&Relative residual\\
\midrule
Radial mass constraint&$2.27319745\,10^{-6}$\\
Lapse constraint&$2.80774843\,10^{-6}$\\
Local mass flux against a time difference&$1.26971871\,10^{-3}$\\
\bottomrule
\end{tabular}
\end{center}
The flux diagnostic is independent of simply inserting the geometry solve's own right side into its output. Its finite time difference is still an approximation. The separately reported absolute outer-mass flux-balance defect is $0.00230311127$. A small total outer-mass change cannot be attributed solely to physical boundary flux unless that flux is integrated and its numerical balance is controlled.

The scalar verification record contains twelve passed checks. Two concern symbolic identities. The remaining checks are bounded numerical tests with declared thresholds. The symbolic audit, null-channel algebra, and seven-row data structure have separate supporting records. Those checks do not constitute a replay of the full collapse or an authentication of its original execution environment. The verification appendix states that evidence boundary.

\section{What remains beyond these strong-field sectors}
\label{c:strong:remaining}

There are exact results here: the electrovac mass identity, its regular-patch constraint equivalence, the reconstructed finite reduced constraints, the scalar flat-slice reduction, the continuum local mass flux, and its compatibility with the radial mass constraint. Their proofs use the displayed chart and symmetry hypotheses. Neither numerical rounding nor a missing fresh Lean build changes them into a machine-checked unrestricted field theory.

There are also controlled numerical observations: a prescribed electrovac gauge trajectory crossing its known outer horizon, and an excised spherical scalar experiment reaching a first marginally trapped sphere. Their retained refinement data and negative control make useful, bounded tests. The next proofs must address the actual boundary closure, regularity margins, vacuum extension, solver conventions, and errors in the full coupled stage map.

For electrovac, a nonvacuous interval energy theorem must control boundary flux and then establish the stronger norm needed by the reconstruction estimate. A full RK4 theorem needs a variable-coefficient argument, rather than the frozen advection symbol alone. For scalar matter, the mass and lapse prefix solves need stability and consistency estimates coupled to the characteristic evolution, including behavior near $v=0$ and the excision boundary. Event convergence needs quantitative control of the compactness maximum and its temporal crossing. A tangential spatial root cannot inherit a simple-root horizon bound automatically.

Beyond these radial sectors, finite-range nonlinear gravitational closure remains a separate problem. Spherical symmetry removes the radiative gravitational degrees of freedom. A general $3+1$ Einstein--Maxwell--Dirac construction must restore them, choose boundary and gauge conditions, establish the relevant constraint algebra and chart regularity, and prove consistency and stability of its actual evolution. Numerical marginal-sphere formation here establishes no global cosmic-censorship result, singularity theorem, nonspherical stability theorem, or unique event-horizon reconstruction.

The exact quasienergy chapter in the source corpus supplies a further boundary of interpretation. Its componentwise momentum change of variables conjugates a branchwise nonlinear Lorentz action to ordinary Lorentz transformations. The resulting momentum-space metric is a pullback of a flat metric. This is a kinematic coordinate construction, not a solution of the gravitational constraints. Likewise, positive curvature of a low-momentum dispersion relation is not spacetime curvature. A gravity coupling of the quasienergy system requires its own tetrad, connection, stress-energy, and dynamical constraint analysis. That obligation is not filled by relabeling a momentum reparametrization as strong-field geometry.

